\chapter{مفاهیم‌ پایه}

در این فصل، مبانی نظری و چارچوب‌های اصلی مورد استفاده در مسئله ابرتفکیک تک‌تصویری مورد بررسی قرار می‌گیرد. هدف این فصل، فراهم‌سازی زمینه‌ای منسجم برای درک روش‌های مورد مطالعه و تحلیل دقیق‌تر آن‌ها در فصول بعدی است.

\section{مدل ریاضی تخریب و بازسازی تصویر}

همان‌گونه که در فصل قبل اشاره شد، مسئله ابرتفکیک را می‌توان به‌صورت معکوس‌سازی یک فرآیند تخریب مدل‌سازی کرد. مدل کلاسیک تخریب به‌صورت زیر بیان می‌شود:

\begin{equation}
	\mathbf{y} = (\mathbf{x} \otimes \mathbf{k}) \downarrow_s + \mathbf{n}
\end{equation}

که در آن $\mathbf{x}$ تصویر با وضوح بالا، $\mathbf{k}$ کرنل تاری، $\downarrow_s$ عملگر کاهش نمونه‌برداری با ضریب $s$ و $\mathbf{n}$ نویز افزوده است. در قالب ماتریسی، این رابطه به شکل زیر بازنویسی می‌شود:

\begin{equation}
	\mathbf{y} = \mathbf{D}\mathbf{x} + \mathbf{n}
\end{equation}

از آنجا که عملگر $\mathbf{D}$ اطلاعات فرکانسی بالا را حذف می‌کند، بازیابی $\mathbf{x}$ از روی $\mathbf{y}$ نیازمند تعریف پیشین مناسب است. انتخاب این پیشین، هسته اصلی تفاوت میان روش‌های مختلف ابرتفکیک را تشکیل می‌دهد.

\section{بازنمایی تنک و یادگیری دیکشنری}

یکی از تأثیرگذارترین چارچوب‌های کلاسیک در ابرتفکیک، مبتنی بر بازنمایی تنک\LTRfootnote{Sparse Representation} است. فرض اساسی در این رویکرد آن است که هر پچ تصویر را می‌توان به‌صورت ترکیب خطی تعداد اندکی اتم از یک دیکشنری بیان کرد.

برای یک پچ وضوح پایین $\mathbf{y}$ و دیکشنری $\mathbf{D}_l$، مسئله بهینه‌سازی به شکل زیر تعریف می‌شود:

\begin{equation}
	\boldsymbol{\alpha}^\star =
	\arg\min_{\boldsymbol{\alpha}}
	\left\|
	\mathbf{y} - \mathbf{D}_l \boldsymbol{\alpha}
	\right\|_2^2
	+
	\lambda
	\left\|
	\boldsymbol{\alpha}
	\right\|_1
\end{equation}

که در آن $\lambda$ پارامتر تنظیم‌کننده تنکی است. فرض کلیدی این روش آن است که ضرایب تنک به‌دست‌آمده در فضای وضوح پایین، در فضای وضوح بالا نیز معتبر هستند:

\begin{equation}
	\mathbf{x} = \mathbf{D}_h \boldsymbol{\alpha}^\star
\end{equation}

که $\mathbf{D}_h$ دیکشنری متناظر در فضای وضوح بالا است.

حل این مسئله معمولاً از طریق الگوریتم‌های بهینه‌سازی محدب یا روش‌های حریصانه انجام می‌شود. پیچیدگی محاسباتی و نیاز به حل مسئله برای تمامی پچ‌های تصویر، از چالش‌های اساسی این چارچوب به شمار می‌آید.

\section{یادگیری عمیق داده‌محور}

با ظهور شبکه‌های کانولوشنی عمیق\LTRfootnote{Convolutional Neural Networks}، مسئله ابرتفکیک به‌صورت یادگیری یک نگاشت غیرخطی از فضای تصاویر وضوح پایین به فضای وضوح بالا مدل شد:

\begin{equation}
	\mathbf{x} = f_\theta(\mathbf{y})
\end{equation}

که در آن $f_\theta$ شبکه‌ای با پارامترهای $\theta$ است که از طریق کمینه‌سازی تابع زیان زیر آموزش داده می‌شود:

\begin{equation}
	\mathcal{L}(\theta) =
	\mathbb{E}
	\left[
	\left\|
	f_\theta(\mathbf{y}) - \mathbf{x}
	\right\|_p
	\right]
\end{equation}

این روش‌ها با استفاده از داده‌های گسترده خارجی آموزش دیده و قادر به مدل‌سازی روابط پیچیده میان الگوهای وضوح پایین و بالا هستند. با این حال، عملکرد آن‌ها به‌شدت وابسته به سازگاری مدل تخریب داده‌های آموزشی با داده‌های آزمون است.

\section{یادگیری درونی و سازگاری تصویر-ویژه}

در مقابل رویکردهای آموزش‌دیده بر داده‌های خارجی، چارچوب یادگیری درونی\LTRfootnote{Internal Learning} بر این اصل استوار است که تصاویر طبیعی دارای بازگشت‌پذیری آماری درونی در مقیاس‌های مختلف هستند. به بیان دیگر، قطعات کوچک تصویر در سطوح مقیاس متفاوت تکرار می‌شوند.

در این چارچوب، مجموعه آموزش مستقیماً از خود تصویر استخراج شده و شبکه‌ای کوچک در زمان آزمون آموزش داده می‌شود. این شبکه نگاشت تصویر-ویژه را می‌آموزد:

\begin{equation}
	\mathbf{x} = f_{\theta_I}(\mathbf{y})
\end{equation}

که در آن پارامترهای $\theta_I$ برای هر تصویر به‌صورت مستقل یاد گرفته می‌شوند. این سازگاری تصویر-محور امکان عملکرد مناسب در شرایط تخریب ناشناخته و غیرایده‌آل را فراهم می‌سازد.

\section{مدل‌سازی تخریب پیچیده و واقع‌گرایی ادراکی}

یکی از محدودیت‌های اساسی مدل‌های کلاسیک، ساده‌سازی بیش‌ازحد فرآیند تخریب است. در تصاویر واقعی، تخریب اغلب شامل چندین مرحله متوالی از تاری، نویز، تغییر اندازه، فشرده‌سازی و پردازش‌های غیرخطی است. این فرآیند را می‌توان به‌صورت زیر مدل کرد:

\begin{equation}
	\mathbf{y} =
	\mathcal{D}_n
	(
	\mathcal{D}_{n-1}
	(
	\dots
	\mathcal{D}_1(\mathbf{x})
	)
	)
\end{equation}

که در آن هر $\mathcal{D}_i$ یک مرحله تخریب با پارامترهای مستقل است. این چارچوب که به مدل تخریب مرتبه‌بالا\LTRfootnote{High-Order Degradation} شناخته می‌شود، فضای وسیع‌تری از شرایط واقعی را پوشش می‌دهد.

برای مقابله با این پیچیدگی، شبکه‌های مولد تخاصمی\LTRfootnote{Generative Adversarial Networks} مورد استفاده قرار گرفتند. در این ساختار، یک شبکه مولد $G$ تصویر ابرتفکیک‌شده را تولید می‌کند و یک شبکه تمایزبخش $D$ کیفیت ادراکی آن را ارزیابی می‌نماید:

\begin{equation}
	\min_G \max_D
	\mathbb{E}_{\mathbf{x}}[\log D(\mathbf{x})]
	+
	\mathbb{E}_{\mathbf{y}}[\log(1 - D(G(\mathbf{y})))]
\end{equation}

نسخه‌های پیشرفته‌تر این چارچوب با طراحی ساختارهای عمیق‌تر برای مولد و تمایزبخش، و همچنین با مدل‌سازی دقیق‌تر فرآیندهای تخریب پیچیده، تلاش کرده‌اند شکاف میان داده‌های مصنوعی و تصاویر واقعی را کاهش دهند. استفاده از ساختارهای باقیمانده درون‌باقیمانده\LTRfootnote{Residual-in-Residual Dense Blocks} در مولد و تمایزبخش‌های مبتنی بر ساختار U شکل با نرمال‌سازی طیفی، نمونه‌ای از این پیشرفت‌ها است.
\section{تعاریف و نمادگذاری‌های متداول در ابرتفکیک}

در ادبیات ابرتفکیک، از مجموعه‌ای از نمادگذاری‌ها و اصطلاحات استاندارد استفاده می‌شود که درک دقیق آن‌ها برای تحلیل روش‌ها ضروری است. در این پژوهش نیز از همین قراردادها پیروی می‌شود.

تصویر با وضوح پایین با نماد $\mathbf{y}$ و عنوان \LTRfootnote{Low-Resolution (LR)} نمایش داده می‌شود. تصویر مرجع با وضوح بالا که به‌عنوان حقیقت مبنا در ارزیابی استفاده می‌شود، با نماد $\mathbf{x}$ یا $\mathbf{x}_{gt}$ و عنوان \LTRfootnote{High-Resolution (HR)} یا \LTRfootnote{Ground Truth (GT)} شناخته می‌شود. خروجی الگوریتم ابرتفکیک با نماد $\hat{\mathbf{x}}$ نمایش داده شده و به آن تصویر \LTRfootnote{Super-Resolved (SR)} گفته می‌شود.

رابطه کلی میان این کمیت‌ها را می‌توان به‌صورت زیر خلاصه کرد:

\begin{equation}
	\hat{\mathbf{x}} = \mathcal{F}(\mathbf{y})
\end{equation}

که در آن $\mathcal{F}$ مدل یا الگوریتم ابرتفکیک است. در ارزیابی کمی، هدف آن است که فاصله میان $\hat{\mathbf{x}}$ و $\mathbf{x}_{gt}$ کمینه شود.

ضریب بزرگ‌نمایی معمولاً با $s$ نمایش داده می‌شود، به‌طوری‌که:

\begin{equation}
	\text{dim}(\mathbf{x}_{gt}) = s \times \text{dim}(\mathbf{y})
\end{equation}

مقادیر رایج برای $s$ برابر ۲، ۳ یا ۴ هستند.

\section{معیارهای ارزیابی کمی}

ارزیابی عملکرد روش‌های ابرتفکیک معمولاً با استفاده از معیارهای کمی مبتنی بر فاصله پیکسلی یا شباهت ساختاری انجام می‌شود.

\subsection{ریشه میانگین مربعات خطا}

یکی از ساده‌ترین معیارها، ریشه میانگین مربعات خطا\LTRfootnote{Root Mean Square Error (RMSE)} است که به‌صورت زیر تعریف می‌شود:

\begin{equation}
	\text{RMSE} =
	\sqrt{
		\frac{1}{N}
		\sum_{i=1}^{N}
		\left(
		\hat{x}_i - x_{gt,i}
		\right)^2
	}
\end{equation}

که در آن $N$ تعداد کل پیکسل‌ها است. این معیار اختلاف مطلق شدت پیکسل‌ها را اندازه‌گیری می‌کند، اما همبستگی آن با ادراک بصری انسان محدود است.

\subsection{نسبت سیگنال به نویز بیشینه}

معیار پرکاربرد دیگر، نسبت سیگنال به نویز بیشینه\LTRfootnote{Peak Signal-to-Noise Ratio (PSNR)} است که بر پایه میانگین مربعات خطا تعریف می‌شود:

\begin{equation}
	\text{PSNR} =
	10 \log_{10}
	\left(
	\frac{L^2}{\text{MSE}}
	\right)
\end{equation}

که در آن:

\begin{equation}
	\text{MSE} =
	\frac{1}{N}
	\sum_{i=1}^{N}
	\left(
	\hat{x}_i - x_{gt,i}
	\right)^2
\end{equation}

و $L$ بیشینه مقدار شدت پیکسل (برای تصاویر ۸ بیتی برابر ۲۵۵) است. مقدار بزرگ‌تر PSNR نشان‌دهنده بازسازی نزدیک‌تر به تصویر مرجع است.

\subsection{شاخص شباهت ساختاری}

به‌منظور لحاظ‌کردن ویژگی‌های ادراکی، شاخص شباهت ساختاری\LTRfootnote{Structural Similarity Index (SSIM)} معرفی شد. این معیار سه مؤلفه روشنایی، کنتراست و ساختار را مقایسه می‌کند:

\begin{equation}
	\text{SSIM}(x,y) =
	\frac{(2\mu_x \mu_y + C_1)(2\sigma_{xy} + C_2)}
	{(\mu_x^2 + \mu_y^2 + C_1)(\sigma_x^2 + \sigma_y^2 + C_2)}
\end{equation}

که در آن $\mu$ میانگین، $\sigma^2$ واریانس و $\sigma_{xy}$ کوواریانس دو تصویر است. SSIM مقداری بین صفر و یک دارد و مقادیر نزدیک به یک نشان‌دهنده شباهت ساختاری بیشتر هستند.

\section{شبکه‌های مولد تخاصمی}

شبکه‌های مولد تخاصمی\LTRfootnote{Generative Adversarial Networks (GANs)} چارچوبی برای یادگیری توزیع داده‌ها از طریق رقابت میان دو شبکه عصبی هستند. این چارچوب شامل یک مولد $G$ و یک تمایزبخش $D$ است.

مولد تلاش می‌کند نمونه‌هایی تولید کند که از نظر آماری مشابه داده‌های واقعی باشند، در حالی که تمایزبخش تلاش می‌کند نمونه‌های واقعی و تولیدی را از یکدیگر تفکیک کند. مسئله بهینه‌سازی به‌صورت زیر تعریف می‌شود:

\begin{equation}
	\min_G \max_D
	\mathbb{E}_{\mathbf{x}\sim p_{data}}
	[\log D(\mathbf{x})]
	+
	\mathbb{E}_{\mathbf{y}}
	[\log(1 - D(G(\mathbf{y})))]
\end{equation}

در ابرتفکیک، مولد نگاشت $\mathbf{y} \rightarrow \hat{\mathbf{x}}$ را می‌آموزد و تمایزبخش کیفیت ادراکی خروجی را ارزیابی می‌کند. برخلاف روش‌های صرفاً مبتنی بر خطای پیکسلی، استفاده از GAN موجب افزایش واقع‌گرایی بافت‌ها و جزئیات ریز می‌شود، هرچند ممکن است مقدار PSNR کاهش یابد.

در نسخه‌های پیشرفته‌تر، علاوه بر زیان تخاصمی، از زیان ادراکی\LTRfootnote{Perceptual Loss} مبتنی بر ویژگی‌های استخراج‌شده از شبکه‌های ازپیش‌آموزش‌دیده نیز استفاده می‌شود:

\begin{equation}
	\mathcal{L}_{perc} =
	\left\|
	\phi(\hat{\mathbf{x}})
	-
	\phi(\mathbf{x}_{gt})
	\right\|_2^2
\end{equation}

که در آن $\phi$ نگاشت ویژگی در یک شبکه عمیق ثابت است.

\section{موازنه دقت عددی و کیفیت ادراکی}

یکی از مباحث مهم در ابرتفکیک، تضاد میان بیشینه‌سازی معیارهای عددی مانند PSNR و بهبود کیفیت ادراکی است. کمینه‌سازی خطای مربعی منجر به خروجی‌های نرم و میانگین‌گرا می‌شود، در حالی که بهینه‌سازی مبتنی بر GAN تمایل به تولید بافت‌های تیزتر اما گاه غیرمنطبق با مقدار دقیق پیکسلی دارد.

این تعادل را می‌توان به‌صورت ترکیب چند تابع زیان مدل کرد:

\begin{equation}
	\mathcal{L}_{total} =
	\lambda_1 \mathcal{L}_{pixel}
	+
	\lambda_2 \mathcal{L}_{perc}
	+
	\lambda_3 \mathcal{L}_{adv}
\end{equation}

که در آن ضرایب $\lambda_i$ وزن هر مؤلفه را تعیین می‌کنند.
\section{ارزیابی بر روی کانال روشنایی در فضای رنگی}

در بسیاری از مقالات ابرتفکیک، ارزیابی کمی به‌جای انجام در فضای رنگی \LTRfootnote{RGB}، بر روی کانال روشنایی در فضای \LTRfootnote{YCbCr} انجام می‌شود. دلیل این انتخاب به نحوه ادراک سیستم بینایی انسان بازمی‌گردد.

فضای YCbCr تصویر را به سه مؤلفه تفکیک می‌کند:
\begin{itemize}
	\item مؤلفه روشنایی $Y$  
	\item مؤلفه اختلاف آبی $Cb$  
	\item مؤلفه اختلاف قرمز $Cr$
\end{itemize}

سیستم بینایی انسان حساسیت بیشتری نسبت به تغییرات روشنایی نسبت به تغییرات رنگی دارد. به همین دلیل، بسیاری از روش‌های ابرتفکیک تنها کانال $Y$ را بازسازی کرده و کانال‌های رنگی را با روش‌های ساده‌تر مانند درونیابی بزرگ‌نمایی می‌کنند.

رابطه تبدیل از RGB به YCbCr به‌صورت خطی بیان می‌شود:

\begin{equation}
	Y = 0.299R + 0.587G + 0.114B
\end{equation}

در ارزیابی کمی، ابتدا تصاویر به فضای YCbCr تبدیل شده و سپس معیارهایی مانند PSNR و SSIM تنها بر روی کانال $Y$ محاسبه می‌شوند:

\begin{equation}
	\text{PSNR}_Y =
	10 \log_{10}
	\left(
	\frac{L^2}{\text{MSE}(Y)}
	\right)
\end{equation}

این رویکرد باعث می‌شود نتایج عددی با ادراک بصری انسان همبستگی بیشتری داشته باشند و همچنین مقایسه عادلانه‌تری میان روش‌هایی که تنها روشنایی را مدل می‌کنند فراهم گردد.

\section{معیارهای ارزیابی بدون مرجع}

معیارهایی مانند PSNR و SSIM نیازمند در اختیار داشتن تصویر مرجع \LTRfootnote{Ground Truth} هستند. با این حال، در بسیاری از کاربردهای واقعی، تصویر وضوح بالا در دسترس نیست. در چنین شرایطی از معیارهای بدون مرجع\LTRfootnote{No-Reference Image Quality Assessment} استفاده می‌شود.

یکی از پرکاربردترین این معیارها، شاخص کیفیت طبیعی تصویر\LTRfootnote{Natural Image Quality Evaluator (NIQE)} است. ایده اصلی NIQE بر این فرض استوار است که تصاویر طبیعی دارای توزیع آماری مشخصی در حوزه ویژگی‌های آماری هستند.

در این روش، ویژگی‌هایی مبتنی بر مدل آماری گوسی تعمیم‌یافته\LTRfootnote{Generalized Gaussian Model} از تصویر استخراج شده و با یک مدل آماری آموزش‌دیده از تصاویر طبیعی مقایسه می‌شوند. فاصله آماری میان این دو توزیع به‌عنوان معیار کیفیت تعریف می‌شود:

\begin{equation}
	\text{NIQE} =
	\sqrt{
		(\mu - \mu_n)^T
		\Sigma^{-1}
		(\mu - \mu_n)
	}
\end{equation}

که در آن $\mu$ و $\Sigma$ میانگین و کوواریانس ویژگی‌های تصویر ارزیابی‌شده و $\mu_n$ و $\Sigma_n$ پارامترهای مدل تصاویر طبیعی هستند.

برخلاف PSNR، در NIQE مقدار کمتر نشان‌دهنده کیفیت بهتر است. این معیار به‌ویژه در ارزیابی روش‌های مبتنی بر شبکه‌های مولد که تمرکز بر کیفیت ادراکی دارند، اهمیت زیادی پیدا می‌کند.

\section{روش بازافکنی سازگار}

یکی از تکنیک‌های کلاسیک برای بهبود سازگاری نتیجه ابرتفکیک با مدل تخریب، روش بازافکنی\LTRfootnote{Back-Projection} است. ایده اصلی این روش آن است که تصویر بازسازی‌شده باید پس از اعمال مدل تخریب، مجدداً تصویر وضوح پایین اولیه را تولید کند.

اگر $\hat{\mathbf{x}}$ خروجی اولیه الگوریتم باشد، می‌توان آن را از طریق مدل تخریب به وضوح پایین بازگرداند:

\begin{equation}
	\mathbf{y}' = \mathbf{D}\hat{\mathbf{x}}
\end{equation}

اختلاف میان $\mathbf{y}'$ و تصویر ورودی $\mathbf{y}$ برابر است با:

\begin{equation}
	\mathbf{e} = \mathbf{y} - \mathbf{y}'
\end{equation}

این خطا سپس به فضای وضوح بالا بازافکنی شده و برای اصلاح تصویر استفاده می‌شود:

\begin{equation}
	\hat{\mathbf{x}}^{(t+1)} =
	\hat{\mathbf{x}}^{(t)}
	+
	\mathbf{D}^T \mathbf{e}
\end{equation}

که $\mathbf{D}^T$ عملگر معکوس تقریبی (ترانهاده) مدل تخریب است. این فرآیند به‌صورت تکراری اجرا می‌شود تا سازگاری میان تصویر SR و تصویر LR اولیه برقرار گردد.

روش بازافکنی به‌ویژه در شرایطی که مدل تخریب شناخته‌شده یا تخمین‌زده‌شده باشد، موجب بهبود دقت عددی و کاهش مصنوعات ناخواسته می‌شود. بسیاری از روش‌های مدرن، حتی در چارچوب یادگیری عمیق، از این ایده به‌صورت یک مرحله اصلاح نهایی بهره می‌برند.

\section{جمع‌بندی}

در این فصل، چارچوب نظری و مفاهیم پایه مورد نیاز برای تحلیل مسئله ابرتفکیک تک‌تصویری به‌صورت منسجم مورد بررسی قرار گرفت. ابتدا مدل ریاضی تخریب تصویر و ماهیت بدوضع این مسئله تشریح شد و نشان داده شد که بازیابی تصویر با وضوح بالا مستلزم بهره‌گیری از پیشین‌های مناسب یا یادگیری نگاشت‌های پیچیده غیرخطی است.

در ادامه، سه جریان اصلی در توسعه روش‌های ابرتفکیک معرفی گردید: رویکردهای مبتنی بر بازنمایی تنک و یادگیری دیکشنری، روش‌های داده‌محور مبتنی بر شبکه‌های کانولوشنی عمیق، و چارچوب‌های تصویر-ویژه مبتنی بر یادگیری درونی. همچنین مدل‌سازی تخریب‌های پیچیده و مرتبه‌بالا و نقش شبکه‌های مولد تخاصمی در بهبود کیفیت ادراکی و بازسازی بافت‌های واقع‌گرایانه مورد بررسی قرار گرفت. نشان داده شد که این رویکردها در تعامل با یکدیگر شکل گرفته‌اند و هر یک بخشی از چالش‌های مسئله را هدف قرار می‌دهند.

علاوه بر مبانی مدل‌سازی، نمادگذاری‌های استاندارد در حوزه ابرتفکیک شامل تصاویر وضوح پایین، وضوح بالا، خروجی بازسازی‌شده و تصویر مرجع معرفی گردید تا چارچوبی دقیق برای تحلیل‌های بعدی فراهم شود. سپس معیارهای ارزیابی کمی مانند RMSE، PSNR و SSIM تشریح شدند و تفاوت آن‌ها از منظر دقت پیکسلی و شباهت ساختاری تبیین گردید. همچنین اهمیت ارزیابی بر روی کانال روشنایی در فضای رنگی YCbCr و دلیل ترجیح آن در ادبیات پژوهش توضیح داده شد.

در ادامه، محدودیت معیارهای مرجع‌دار در سناریوهای واقعی مطرح و شاخص‌های بدون مرجع نظیر NIQE معرفی شدند که امکان ارزیابی کیفیت ادراکی در غیاب تصویر حقیقت مبنا را فراهم می‌کنند. همچنین مفهوم بازافکنی سازگار به‌عنوان مکانیزمی برای تضمین انطباق خروجی ابرتفکیک با مدل تخریب اولیه بررسی شد.

برآیند مباحث این فصل نشان می‌دهد که مسئله ابرتفکیک تنها یک مسئله بازسازی عددی ساده نیست، بلکه تعادلی میان مدل‌سازی دقیق تخریب، بهره‌گیری از پیشین‌های آماری یا ساختاری، یادگیری نگاشت‌های پیچیده و ارزیابی چندبعدی کیفیت را طلب می‌کند. این چارچوب نظری زمینه لازم برای بررسی، تحلیل و بهبود روش‌های منتخب در فصول آتی را فراهم می‌سازد.